\section{Prueba del Teorema~\ref{teo:normalidad}}\label{app:prueba}

El estimador doblemente robusto es $\theta$ más una media de términos i.i.d. de media cero, más restos de orden $o_P(n^{-1/2})$; el teorema central del límite hace el resto.

\paragraph{Notación.} Trabajamos primero con observaciones i.i.d. $O=(Z,X,(1-Z)Y)$; el caso con varias observaciones por unidad se trata al final. $P_nf=n^{-1}\sum_if(O_i)$; $\E f$ es la esperanza poblacional; $\|h\|_Q$ y $\|h\|_P$ son normas de $L^2(Q)$ y $L^2(P)$. Sean $o_{\max}=(1-\eta)/\eta$ y $w_{\max}=(1-\pi)o_{\max}/\pi$: por la condición 1 y el Teorema~\ref{teo:odds}, $o\le o_{\max}$ y $w\le w_{\max}$, y por la condición 2, $\hat o\le o_{\max}$. Las observaciones están acotadas: $|Y|\le B_Y$ y $|\hat m|\le B_Y$. En el bloque $k$, con $n_k$ observaciones, $(\hat o_k,\hat m_k)$ se entrenan con los datos de fuera del bloque; $\E_k$ y $\PP_k$ denotan esperanza y probabilidad condicionadas a esos datos, y $P_{n,k}$ la media empírica dentro del bloque.

\paragraph{Tres identidades.} Para toda $h$ medible y acotada:
\begin{align*}
&\text{(I1)}\quad \E[(1-Z)\,o(X)\,h(X)]=\pi\,\E_P[h],\\
&\text{(I2)}\quad \E[(1-Z)\,h(X)\,(Y-m_0(X))]=0,\\
&\text{(I3)}\quad \|h\|_P^2\le w_{\max}\,\|h\|_Q^2 .
\end{align*}
(I1): $\E[(1-Z)oh]=(1-\pi)\E_Q[oh]$, y $o=\pi w/(1-\pi)$ por el Teorema~\ref{teo:odds}, luego vale $\pi\E_Q[wh]=\pi\E_P[h]$. (I2): torre condicionando en $X$ sobre $\{Z=0\}$, donde $\E[Y\mid X,Z=0]=m_0(X)$. (I3): $\|h\|_P^2=\E_Q[wh^2]\le w_{\max}\,\E_Q[h^2]$.

\paragraph{Paso 1 (la función $\psi$ y su sesgo).} Para funciones acotadas $(o',m')$, sea $\psi(O;o',m')=(1-Z)\,o'(X)\,(Y-m'(X))+Z\,(m'(X)-\theta)$. Entonces $\pi\phi=\psi(\cdot;o,m_0)$ y
\[
\E[\psi(O;o',m')]=(1-\pi)\,\E_Q[(o'-o)(m_0-m')].
\]
Por (I2) con $h=o'$, $\E[(1-Z)o'(Y-m')]=(1-\pi)\E_Q[o'(m_0-m')]$. Por (A2), $\theta=\E_P[m_0]$, y por (I1) con $h=m'-m_0$, $\E[Z(m'-\theta)]=\pi\E_P[m'-m_0]=(1-\pi)\E_Q[o(m'-m_0)]$. Sumando se obtiene la fórmula; en particular, $\E\phi=0$. Es la Proposición~\ref{prop:dr} escrita con odds.

\paragraph{Paso 2 (linealización).} Con $\hat A=P_n[(1-Z)\hat o(Y-\hat m)]$, $\hat B=P_n[(1-Z)\hat o]$ y $\hat\pi=P_nZ$, el estimador es $\hat\theta=\hat A/\hat B+P_n[Z\hat m]/\hat\pi$. En el suceso $\{\hat\pi>\pi/2,\ \hat B>\pi/2\}$, cuya probabilidad tiende a 1 y en el que trabajamos en adelante, $P_n[Z\theta]/\hat\pi=\theta$ y el álgebra da exactamente
\[
\hat\theta-\theta=\frac1\pi P_n\psi(O;\hat o,\hat m)+\hat A\Bigl(\frac1{\hat B}-\frac1\pi\Bigr)+P_n[Z(\hat m-\theta)]\Bigl(\frac1{\hat\pi}-\frac1\pi\Bigr).
\]
Los pasos 3 y 4 muestran que el primer término es $P_n\phi+o_P(n^{-1/2})$; el paso 5, que los otros dos son $o_P(n^{-1/2})$.

\paragraph{Paso 3 (proceso empírico).} Sea $\Delta_k=\psi(\cdot;\hat o_k,\hat m_k)-\psi(\cdot;o,m_0)$. Condicionado a los datos de fuera del bloque $k$, $\hat o_k$ y $\hat m_k$ son funciones fijas y los $\Delta_k(O_i)$ del bloque son i.i.d. Por la desigualdad de Chebyshev condicional, para todo $\epsilon>0$,
\[
\PP_k\Bigl(\bigl|(P_{n,k}-\E_k)\Delta_k\bigr|>\epsilon\,n^{-1/2}\Bigr)\le\frac{n}{n_k}\,\frac{\E_k[\Delta_k^2]}{\epsilon^2}.
\]
Como $\Delta=(1-Z)[(\hat o-o)(Y-m_0)+\hat o(m_0-\hat m)]+Z(\hat m-m_0)$ y $|Y-m_0|\le2B_Y$, la desigualdad $(a+b+c)^2\le3(a^2+b^2+c^2)$ y (I3) dan
\[
\E_k[\Delta_k^2]\le3\Bigl[4B_Y^2(1-\pi)\,\|\hat o_k-o\|_Q^2+\bigl((1-\pi)o_{\max}^{\,2}+\pi w_{\max}\bigr)\|\hat m_k-m_0\|_Q^2\Bigr],
\]
que tiende a 0 en probabilidad por la condición 3. Como $n/n_k\to K$, la probabilidad condicional tiende a 0 en probabilidad, y al estar acotada por 1, la convergencia dominada da que la probabilidad sin condicionar también tiende a 0. Luego $(P_{n,k}-\E_k)\Delta_k=o_P(n^{-1/2})$ para cada uno de los $K$ bloques.

\paragraph{Paso 4 (sesgo).} Como $\E_k\psi(\cdot;o,m_0)=0$, el paso 1 y Cauchy--Schwarz dan
\[
|\E_k\Delta_k|=(1-\pi)\bigl|\E_Q[(\hat o_k-o)(m_0-\hat m_k)]\bigr|\le(1-\pi)\,r_o\,r_m=o_P(n^{-1/2}).
\]
Sumando sobre los bloques con pesos $n_k/n$, $P_n\psi(\cdot;\hat o,\hat m)=P_n\psi(\cdot;o,m_0)+o_P(n^{-1/2})=\pi P_n\phi+o_P(n^{-1/2})$.

\paragraph{Paso 5 (normalizaciones).} Descomponemos los factores del paso 2 en una parte i.i.d. y una de error de los parámetros auxiliares:
\begin{itemize}
  \item $\hat A=P_n[(1-Z)o(Y-m_0)]+P_n[(1-Z)\{(\hat o-o)(Y-m_0)+\hat o(m_0-\hat m)\}]$. El primer sumando es $O_P(n^{-1/2})$ por el teorema central del límite. En el segundo, la parte centrada es $o_P(n^{-1/2})$ con el argumento del paso 3, y su media condicional, por (I2), es $(1-\pi)\E_Q[\hat o_k(m_0-\hat m_k)]$, acotada por $(1-\pi)o_{\max}\,r_m$. Luego $\hat A=O_P(n^{-1/2})+O_P(r_m)$.
  \item Del mismo modo, $\hat B-\pi=(P_n-\E)[(1-Z)o]+P_n[(1-Z)(\hat o-o)]=O_P(n^{-1/2})+O_P(r_o)$, usando $\E[(1-Z)o]=\pi$ por (I1) y que la norma de $L^1(Q)$ no supera a la de $L^2(Q)$.
  \item $P_n[Z(\hat m-\theta)]=P_n[Z(m_0-\theta)]+P_n[Z(\hat m-m_0)]=O_P(n^{-1/2})+O_P(r_m)$, usando $\E[Z(m_0-\theta)]=0$ y (I3). Además, $1/\hat\pi-1/\pi=O_P(n^{-1/2})$.
\end{itemize}
Como $\hat B\to\pi>0$ en probabilidad,
\[
\hat A\Bigl(\frac1{\hat B}-\frac1\pi\Bigr)=\frac{\hat A(\pi-\hat B)}{\pi\hat B}=O_P\bigl(n^{-1}+n^{-1/2}(r_o+r_m)+r_or_m\bigr)=o_P(n^{-1/2}),
\]
y el tercer término del paso 2 es $O_P(n^{-1/2})\cdot[O_P(n^{-1/2})+O_P(r_m)]=o_P(n^{-1/2})$.

\paragraph{Paso 6 (conclusión).} $\hat\theta-\theta=P_n\phi+o_P(n^{-1/2})$. Las $\phi(O_i)$ son i.i.d., con $\E\phi=0$ y $\E\phi^2=\mathcal V<\infty$, porque $\phi$ está acotada. El teorema central del límite y el lema de Slutsky dan $\sqrt n(\hat\theta-\theta)\to N(0,\mathcal V)$.

\paragraph{Paso 7 (estimación de $\mathcal V$).} El estimador implementado usa
\[
\hat\phi(O)=\frac{(1-Z)\,\hat o(X)\,\bigl(Y-\hat m(X)-\hat A/\hat B\bigr)}{\hat B}+\frac{Z\,(\hat m(X)-\hat\mu_1)}{\hat\pi},\qquad\hat\mu_1=\frac{P_n[Z\hat m]}{\hat\pi}.
\]
Todas las funciones están acotadas, así que $|\hat\phi-\phi|\le C_0\,(|\hat o-o|+|\hat m-m_0|+|\hat\mu_1-\theta|+|\hat A/\hat B|+|\hat B-\pi|+|\hat\pi-\pi|)$ para una constante $C_0$ (aleatoria, pero acotada en ese suceso). Las medias condicionales de los cuadrados de los dos primeros términos tienden a 0 por bloques, y los escalares tienden a 0 en probabilidad, así que $P_n(\hat\phi-\phi)^2=o_P(1)$. Con $P_n\phi^2\to \mathcal V$ por la ley de los grandes números y Cauchy--Schwarz, $|P_n\hat\phi^2-P_n\phi^2|\le P_n(\hat\phi-\phi)^2+2[P_n(\hat\phi-\phi)^2\cdot P_n\phi^2]^{1/2}\to0$. Así $\hat{\mathcal V}\to \mathcal V$ en probabilidad, y por Slutsky el intervalo $\hat\theta\pm z_{1-\alpha/2}(\hat{\mathcal V}/n)^{1/2}$ tiene cobertura asintótica $1-\alpha$.

\paragraph{Varias observaciones por unidad.} Si cada unidad $i$ aporta un bloque de $N_i\le N_{\max}$ observaciones con el mismo $Z_i$, las unidades son i.i.d. y $P$ y $Q$ se definen como leyes a nivel de observación, todo lo anterior se repite con cantidades de bloque. Las identidades (I1)--(I3) valen para sumas de bloque por linealidad, con $\pi$ y $1-\pi$ sustituidos por $\bar N_1=\E[\sum_jZ_{ij}]$ y $\bar N_0=\E[\sum_j(1-Z_{ij})]$, los números esperados de observaciones objetivo y donantes por unidad (que pueden superar 1), y $o$ por los odds a nivel de observación, de modo que $o=\bar N_1w/\bar N_0$. El \emph{cross-fitting} por unidad mantiene la independencia entre el bloque evaluado y los datos de entrenamiento, y $N_i\le N_{\max}$ mantiene todo acotado. La función de influencia por unidad es $\sum_j\psi(O_{ij};o,m_0)/\bar N_1$, y $\sqrt n$ se refiere al número de unidades; el error típico agrupado es su versión muestral. Con dependencia dentro de las unidades, $\mathcal V$ es la varianza de esta función de influencia por bloques, no la forma cerrada de la sección~\ref{sec:inferencia}, que corresponde al caso i.i.d.
